\appendix
\setcounter{table}{0}\renewcommand{\thetable}{A\arabic{table}}

\section{Supplementary tables}\label{app:tables}

\begin{table}[h]
\centering\small
\caption{$B$-parameters at the benchmark ($\theta=6.1\deg$, $\mFP/H_0=25$, $\alpha=0$, $\beta=3$, $\OL=0.699$, $\Orad=9.08\times10^{-5}$).}\label{tab:A1}
\begin{tabular}{@{}lrl@{}}
\toprule
Parameter & Value & Formula\\
\midrule
$B_0$ & $-40.248$ & $3\OL-\sin^2\!\theta\,\mFP^2(3+3\alpha+\beta)$\\
$B_1$ & $28.230$ & $\sin^2\!\theta\,\mFP^2(1+2\alpha+\beta)$\\
$B_2$ & $-21.173$ & $-\sin^2\!\theta\,\mFP^2(\alpha+\beta)$\\
$B_3$ & $21.173$ & $\sin^2\!\theta\,\mFP^2\beta$\\
$B_4$ & $-28.206$ & $3\tan^2\!\theta\,\OL+\sin^2\!\theta\,\mFP^2(-1+\alpha-\beta)$\\
$B_1^2$ & $796.94$ & \\
\bottomrule
\end{tabular}
\end{table}

\begin{table}[h]
\centering\small
\caption{Rigidity of the anchor-implied $\Om$ over the $(\alpha,\beta)$ grid: $\Om$ from the quartic closure at $z=0$ ($\theta=6.1\deg$; CMB convention, neutrinos excluded from $\Om$, reference $h=0.678$). The cell $(\alpha,\beta)=(-1,1)$ carries no finite branch (the degenerate locus $1+2\alpha+\beta=0$, $B_1=0$) and is marked ``--'' rather than filled with a number, consistently with \cref{sec:baores}.}\label{tab:A1b}
\begin{tabular}{@{}rcccc@{}}
\toprule
$\alpha\backslash\beta$ & $1$ & $2$ & $3$ & $4$\\
\midrule
$-1.0$ & -- & $0.3029$ & $0.3029$ & $0.3029$\\
$-0.5$ & $0.3029$ & $0.3029$ & $0.3029$ & $0.3029$\\
$0.0$ & $0.3029$ & $0.3029$ & $0.3029$ & $0.3029$\\
$+0.5$ & $0.3029$ & $0.3029$ & $0.3029$ & $0.3029$\\
$+1.0$ & $0.3029$ & $0.3029$ & $0.3029$ & $0.3029$\\
\bottomrule
\end{tabular}
\end{table}
The $(\alpha,\beta)$-dependent terms in $\ODE(0)$ are of order $(1-y(0))^2=2\times10^{-7}$. No choice of potential parameters within this class moves the anchor-implied $\Om$ towards the SNe-preferred $0.331$ ($z_{\rm HD}$, full covariance).

\begin{table}[h]
\centering\small
\caption{$y(z)$ and $E^2(z)$: bimetric against \LCDM{}. The fractional correction is $\delta E^2=(E^2_{\rm bim}-E^2_{\Lambda{\rm CDM}})/E^2_{\Lambda{\rm CDM}}$ with the same radiation density ($\Orad=9.08\times10^{-5}$) and the same matter density ($\Om=0.3043547$) in both columns. Each column is closed on itself: the \LCDM{} column carries $\OL=1-\Om-\Orad=0.6955545$, the bimetric column $\OL=0.699$.}\label{tab:A2}
\begin{tabular}{@{}rcccc@{}}
\toprule
$z$ & $y(z)$ & $E^2_{\rm bim}$ & $E^2_{\Lambda{\rm CDM}}$ & $\delta E^2\ [\%]$\\
\midrule
$0.00$ & $0.999512$ & $1.000000$ & $1.000000$ & $-0.0000$\\
$0.10$ & $0.999350$ & $1.099643$ & $1.100784$ & $-0.1036$\\
$0.50$ & $0.998352$ & $1.715027$ & $1.723211$ & $-0.4750$\\
$1.00$ & $0.996093$ & $3.107718$ & $3.131845$ & $-0.7704$\\
$2.00$ & $0.986814$ & $8.830855$ & $8.920487$ & $-1.0048$\\
$2.33$ & $0.981968$ & $11.821478$ & $11.945333$ & $-1.0368$\\
$3.00$ & $0.968770$ & $19.980325$ & $20.197501$ & $-1.0753$\\
$5.00$ & $0.896071$ & $65.815888$ & $66.553851$ & $-1.1088$\\
$10.00$ & $0.556527$ & $403.379176$ & $407.121089$ & $-0.9191$\\
$12.00$ & $0.442652$ & $666.804264$ & $671.956212$ & $-0.7667$\\
\bottomrule
\end{tabular}
\end{table}
The column difference is $E^2_{\rm bim}-E^2_{\Lambda{\rm CDM}}=+0.0034455-7.0575\,G(\delta(z))$: the $\Lambda$ offset of $+3.4\times10^{-3}$ is cancelled at $z=0$ by the interaction term through the closure $E^2(0)=1$ (net zero), the difference is $-0.74$ at $z=5$, and tends to $-14.11$ asymptotically (the interaction term $\sin^2\!\theta\,\mFP^2\,G(\delta)$ tends to $2\sin^2\!\theta\,\mFP^2=14.115$ as $\delta\to1$, overwhelming the $+0.0034455$ offset). The ``same radiation density'' clause alone does not fix the comparison, and both pieces of the offset are disclosed here. The maximum magnitude, $1.11\%$, is reached near $z=5$; within the BAO/SNe redshift ranges the correction is at most $1.04\%$. The correction is negative (the bimetric expansion is slower) at all $z>0$.

\begin{table}[h]
\centering\small
\caption{Effective dark-energy equation of state $w_{\rm DE}(z)$. The phantom regime ($w_{\rm DE}<-1$) holds on the positive-$\ODE$ branch up to the zero crossing; $\ODE\to0$ at $z=4.882$, producing a pole in $w_{\rm DE}$; beyond it $\ODE<0$, so the conventional dark-energy interpretation is no longer appropriate.}\label{tab:A3}
\begin{tabular}{@{}rrrrrrr@{}}
\toprule
$z$ & $0.0$ & $0.5$ & $1.0$ & $2.0$ & $3.0$ & $4.0$\\
\midrule
$w_{\rm DE}(z)$ & $-1.0050$ & $-1.0169$ & $-1.0411$ & $-1.1536$ & $-1.4610$ & $-2.5988$\\
\bottomrule
\end{tabular}
\end{table}

\begin{table}[h]
\centering\small
\caption{Accumulation of $D_M(z_{\rm rec})$ by redshift: fraction of the comoving distance contributed by $z\le z_{\rm check}$. The unscreened region ($z\lesssim12$, where $H<\mFP$) contributes $72.2\%$ of $D_M(z_{\rm rec})$ while spanning $\sim1\%$ of the redshift range.}\label{tab:A4}
\begin{tabular}{@{}rrl@{}}
\toprule
$z\le z_{\rm check}$ & $\%$ of $D_M(z_{\rm rec})$ & Regime\\
\midrule
$1$ & $24.32$ & unscreened\\
$5$ & $57.19$ & unscreened\\
$10$ & $69.42$ & unscreened\\
$12$ & $72.19$ & screening boundary\\
$50$ & $87.85$ & screened\\
$500$ & $98.51$ & screened\\
$1090$ & $100.00$ & $z_{\rm rec}$ (last scattering)\\
\bottomrule
\end{tabular}
\end{table}

\begin{table}[h]
\centering\small
\caption{CMB $\theta$ scan: $\chi^2$ (equal to $\Dchi$ against the saturated free reference) at each angle under the adopted $\lambda=0$ protocol. $h$ is the best-fit value; $\Om$ is the CMB closure (neutrinos excluded from $\Om$ and added as $\omega_\nu/h^2$; $\Orad(h)=\omega_r/h^2$) at the fitted $h$ of the row, its growth with $\theta$ being the $\sin^2\!\theta\,\mFP^2G(y(0))$ dependence of the closure (values regenerated by \code{src/checks.py}). The $\lambda^*$ entries are local-parabola estimates from the five-point sweep at fixed $\theta$ (\code{results/cmb\_lambda\_theta\_sweep.json}). The quoted $0.0086$ at $10\deg$ is a parabolic residual; values below $\sim2.5$ at higher angles are the best of five sampled $\lambda$ values, not continuous certificates. The sweep supports structural saturation within the compressed diagnostic, not precise continuous vertices.}\label{tab:A5}
\begin{tabular}{@{}rrccl@{}}
\toprule
$\theta\ [\deg]$ & $\chi^2_{\rm bim}$ & $h$ & $\Om$ & $\lambda^*$\\
\midrule
$6.1$ & $0.00017$ & $0.68688$ & $0.30303$ & $<0$ (boundary; data prefer a higher density)\\
$10.0$ & $10.66599$ & $0.68925$ & $0.30897$ & $0.48$\\
$12.2$ & $34.17418$ & $0.69116$ & $0.31370$ & $0.58$\\
$15.0$ & $98.43996$ & $0.69430$ & $0.32128$ & $0.65$\\
$18.3$ & $253.40524$ & $0.69915$ & $0.33264$ & $0.69$\\
$20.0$ & $381.75291$ & $0.70222$ & $0.33962$ & $0.71$\\
$25.0$ & $1052.78018$ & $0.71413$ & $0.36532$ & $0.74$\\
$30.0$ & $2412.98481$ & $0.73206$ & $0.40039$ & $0.75$\\
\bottomrule
\end{tabular}
\end{table}

\section{Pseudocode for the \texorpdfstring{$E^2(0)$}{E2(0)} residual computation}\label{app:pseudo}
The quartic at redshift $z$, with $t=\tan^2\theta$, $m^2=\mFP^2$, $s^2=\sin^2\theta$ and $\Omega=\Om(1+z)^3+\Orad(1+z)^4$ \citep[][Eq.~5.5]{Hogas2021}, reads
\begin{equation}
\begin{aligned}
c_4&=t\,s^2m^2\beta, & c_3&=3tB_2-B_4, & c_2&=3(tB_1-B_3),\\
c_1&=t(B_0+3\Omega)-3B_2, & c_0&=-B_1,
\end{aligned}
\end{equation}
with $(B_0,\dots,B_4)$ from \citet[][Eq.~3.4a--e]{Hogas2021}. At the benchmark ($t=0.0114210$) the numerical values are $c_4=0.241813$, $c_3=27.480744$, $c_2=-62.550547$, $c_1=63.068553+0.034263\,(\Omega-\Omega(0))$, $c_0=-28.2301$, with $\Omega(0)=0.3044455$ the full non-DE density at $z=0$. These constants are regenerated by \code{src/checks.py}, which reconstructs $E^2(0)$ from them and asserts the closure.
\begin{enumerate}
\item Input: $\theta$, $\mFP$, $\alpha$, $\beta$, $\OL$, $\Orad$.
\item Compute the $B$-parameters (\citealp{Hogas2021}, Eq.~3.4a--e).
\item Find the self-consistent $\Om$: solve $E^2(0)=1$ by Brent's method, evaluating $E^2(0)$ through the quartic root $y(0)$.
\item Residual: $|E^2(0)-1|$; expected below $10^{-12}$.
\item Round trip: recover $(\mFP^2,\OL,\alpha,\beta)$ from the $B$-parameters; agreement expected below $10^{-10}$.
\end{enumerate}

\section{Reproducibility}\label{app:repro}
The code base \code{bimetric\_r28} reproduces every number in this paper from the public data files through documented staged commands, with per-stage return codes, SHA-256 provenance and resumable checkpoints rather than a single monolithic run. \code{src/checks.py} enforces 84 internal checks, \code{src/extra\_numbers.py} adds 69 further checks. The figure generator verifies 117 plotted-number assertions, and the self-similar reduction has 17 dedicated checks. Together these gates cover closure and the B-parameter round trip, probe-level references, scalar values, a Planck regression test, self-similar constants and exact $\varepsilon$-linearity, the ghost-free scan and certificate, plateau flatness, the $\ln(1+z_*)$ scaling and the Monte-Carlo anchor calibration. An independent background implementation (\code{src/independent\_bg.py}) agrees with the canonical solver to $4\times10^{-15}$. Machine-readable assertions accompany the code. The solver raises on failure everywhere; there is no silent fallback. The pipeline is deterministic (fixed grids and seeds, one core per stage).

\section{Self-similar reduction: reduced quartic, lemmas, checks}\label{app:selfsim}
This appendix records only the benchmark specialization used by the pipeline, with exact algebraic/root certificates where stated and numerical cross-checks (17 dedicated checks); arbitrary interaction parameters are not claimed. The reduced problem works directly with the coefficients \eqref{eq:quartic-coeffs} and is independent of the full quartic solver except for the cross-checks. The three lemmas, in the form the checks enforce them, are as follows.

\textbf{Lemma 1 (exact reduction).} After division by $\sin^2\!\theta\,\mFP^2$ the quartic has the coefficients
\begin{equation}
\begin{aligned}
c_4&=t\beta,\qquad c_0=-(1+2\alpha+\beta),\qquad
c_2=3\bigl[t(1+2\alpha+\beta)-\beta\bigr],\\
c_3&=-3t(\alpha+\beta)+(1-\alpha+\beta)-3t\varepsilon,\\
c_1&=t\bigl[-(3+3\alpha+\beta)+3u\bigr]+3(\alpha+\beta)+3t\varepsilon,
\end{aligned}\label{eq:quartic-coeffs}
\end{equation}
and the $\varepsilon$-dependence is exactly linear,
\begin{equation}
P(y;u,\varepsilon)=P(y;u,0)+3t\varepsilon\,(y-y^3),\label{eq:Peps}
\end{equation}
with no $O(\varepsilon^2)$ remainder. The universal root $y(u)$ is obtained by Brent's method on \eqref{eq:quartic-coeffs} at $\varepsilon=0$; $\dd y/\dd u=-3ty/(\partial_yP)$ and $\dd y/\dd N=-3u\,\dd y/\dd u=9tuy/(\partial_yP)$ follow by implicit differentiation (the factor $y$ in $\partial P/\partial u$ is mandatory). Gates: for three test $\varepsilon$ the coefficient shifts are $\Delta c_3=-3t\varepsilon$, $\Delta c_1=+3t\varepsilon$, all others zero to machine precision; $\dd y/\dd N$ from this formula matches finite differences of $y(u)$ to $4\times10^{-11}$.

\textbf{Algebraic ghost-free certificate (benchmark).} For $t=\tan^2(6.1\deg)$ define
\begin{equation*}
\begin{aligned}
F_2&=3y^2-6y+4,\qquad g=8y^3-9y^2+4,\qquad F_3=3ty^2F_2+g,\\
F_4&=4+(6t-5)y+(4-6t)y^2+3ty^3 .
\end{aligned}
\end{equation*}
On $0<y<1$, $F_2=3(y-1)^2+1>0$ and $g\ge37/16>0$, hence $F_3>0$ for $t>0$. Writing $F_4=Q(y)+3ty^3$ with $Q=4+(6t-5)y+(4-6t)y^2$, an exact rational enclosure of $t$ certifies $Q>0$: its interior-vertex minimum is positive and increases for $0<t<1/2$. Thus $F_4>0$. The reduced $\varepsilon=0$ quartic gives $N_4=3tuy-P=(1-y)F_4$, so $\dd y/\dd N=3yN_4/F_3\ge0$ on the open benchmark branch. All sign decisions use exact rational arithmetic in \code{src/certify\_ghost\_full\_branch.py}; the full-model numerical scan is a separate cross-check, not a claim of exact full-model proof. This reduced certificate does not assert $X\ge0$ or full viability.

\textbf{Lemma 2 (boundary).} The reduced universal $X_{\rm red}(u)$ is positive as $u\to0$, crosses zero once at $u_*=30.337301$ (benchmark branch, $\varepsilon=0$), stays negative thereafter, and tends to the analytic asymptote $-1$ (dominant balance $y\sim(1+2\alpha+\beta)/(3tu)$: $T_1\to1$, $T_2\to3$; verified at $u=10^8$ to $3.5\times10^{-6}$). Numerical monotonicity on the scanned branch is a cross-check, not the uniqueness argument. The boundary function is rational in $y$: solving the quartic for $u$ (it enters $c_1$ linearly),
\begin{equation*}
u(y)=\frac{(1-y)\bigl(3ty^3-6ty^2+6ty+4y^2-5y+4\bigr)}{3ty},
\end{equation*}
and the frequency reduces exactly (symbolically verified to $10^{-14}$) to
\begin{equation*}
X(y)=-\frac{Q(y)}{\bigl(3y^2-6y+4\bigr)A(y)^2},\qquad
A(y)=9ty^4-18ty^3+12ty^2+8y^3-9y^2+4,
\end{equation*}
where $Q$ is the following degree-10 polynomial in $(t,y)$:
\begin{equation*}
\begin{aligned}
Q(y)={}&64+(-288+192t)y+(744-912t)y^2+(-1376+2304t)y^3\\
&+(1746-4140t+288t^2)y^4+(-1542+5340t-1296t^2)y^5\\
&+(979-4842t+2376t^2)y^6+(-432+2976t-2376t^2)y^7\\
&+(96-1152t+1404t^2)y^8+(216t-486t^2)y^9+81t^2y^{10}.
\end{aligned}
\end{equation*}
The same coefficient list is machine-readable as \code{src/selfsim.py:q\_coeffs} (stored there in descending-power order). Both denominator factors are strictly positive on $0<y<1$, so the zeros of $X$ on the branch are exactly the roots of $Q$. At the benchmark $t$, exact rational Sturm sequences give one root in $(0,1)$, independently confirmed by \code{sympy.count\_roots}. Moreover, this is an exact interval result rather than only a check at a chosen rational representative: $Q(0)=64$, $Q(1)=-9(t+1)^2\neq0$ on the whole certified interval, and the discriminant of $Q(y)$ with respect to $y$ is nonzero. Therefore the adopted decimal $t=\tan^2(6.1^\circ)$ is covered by an exact interval uniqueness certificate for the reduced polynomial. The numerical estimate is $y_*\simeq0.705464788$. The branch map $y\leftrightarrow u$ is one-to-one: the algebraic certificate gives $\dd y/\dd N\ge0$, strictly positive on $0<y<1$. Combining this bijection with the positive denominator, exactly one Sturm root, and the endpoint signs $X>0$ as $u\to0$ and $X\to-1$ as $u\to\infty$ certifies existence and uniqueness of the crossing. Numerical monotonicity and full-model agreement are cross-checks. The full-model $X$ at finite $z$ carries the additive radiation correction $\rho(z)$ (so the numerical monotonicity statement is made for $X_{\rm red}$ only). Polynomial boundary conditions for restricted submodels exist in the literature (a quintic in $y$ for the $\beta_1$-$\beta_2$ model, \citealp[Eq.~39]{Akrami2015}; closed algebraic eigenfrequencies with a Descartes count, \citealp{Konnig2014}); the rational representative for the full model (all $B_n\neq0$) and its Sturm count are new here. The full model at $\mFP/H_0=2000$, $\Orad=0$, reproduces $u_*$ to $10^{-5}$ and the plateau ratio $H(z_*)/\mFP$ to $5\times10^{-8}$.

\textbf{Lemma 3 (constants).} With $F(\varepsilon)=u_*(\varepsilon)-G(\delta_*(\varepsilon))$ the Friedmann equation at the boundary reads $H^2(z_*)=\sin^2\!\theta\,\mFP^2F(\varepsilon)+\OL$; expanding $F$ to first order at $\varepsilon=0$ at fixed $\OL$ gives \eqref{eq:Hlaw} with $\kappa=1+F'(0)$. Gates (one line): $c=0.582199$, $\kappa=0.374329$ (central differences, step $10^{-5}$ in $\varepsilon$; the decomposition $1-0.626809+0.001138$ is recovered term by term); the plateau is flat to $9\times10^{-7}$ over $\mFP/H_0=500$--$8000$ at $\Orad=0$; the scaling slope of $\ln(1+z_*)$ against $\ln\mFP$ is $0.6669\pm0.0025$ against $2/3$ (the $0.6665$ quoted in \cref{sec:selfsim} is the independent measurement of the \cref{sec:growth} scan); \eqref{eq:Hlaw} reproduces the full-model $H(z_*)$ at the benchmark to $+4.6\times10^{-4}$ relative; the master equation \eqref{eq:master} reconstructs $z_*$ from the reduced problem alone across $\mFP/H_0=3$--$200$ with a residual of $0.49\%$ at most; and the ghost-free check scans $\dd y/\dd N$ over $0\le z\le1090$ on the full branch: $\dd y/\dd N(0)=1.4648\times10^{-3}$, minimum $5.104\times10^{-6}$ at $z=1090$. Applicability: the asymptotic H-law and first-order master expansion assume $\varepsilon\ll1$, $\rho\ll1$. Practical reduced-model tests extend to $\varepsilon\lesssim1$; approximate numerical linearity of $u_*(\varepsilon)$ persists to $\varepsilon\sim7$, without enlarging the asymptotic domain. The radiation term reaches $\rho(1090)\sim0.3$. Negative result, enforced by construction: there is no $I_1$ check -- the $\OL$ and radiation terms of the growth integrand do not reduce to functions of $u$.
